\chapter{Wielandt Bound}\label{ch:wielandt}

This chapter proves a quantum analog of Wielandt's inequality:
for a normal MPS tensor with bond dimension~$D$,
the cumulative span of matrix products stabilises
after at most $D^2$ steps.
The results follow~\cite{Sanz2010Quantum}.

\section{Cumulative span}

\begin{definition}[Word span]\label{def:word_span}
    \lean{MPSTensor.wordSpan}
    \leanok
    \uses{def:mps_tensor, def:eval_word}
    The \emph{word span} at length~$n$ is the subspace
    \[
        S_n(A) =  \spn_{\C}\{A^w : |w| = n\}
        \subseteq  \MN{D}.
    \]
    This is \cite[Eq.~(1)]{Sanz2010Quantum}.
\end{definition}

\begin{definition}[Cumulative span]\label{def:cumulative_span}
    \lean{MPSTensor.cumulativeSpan}
    \leanok
    \uses{def:word_span}
    The \emph{cumulative span} at level~$n$ is
    \[
        T_n(A) =  S_0(A) + S_1(A) + \cdots + S_n(A).
    \]
\end{definition}

\begin{lemma}[Monotonicity of cumulative span]\label{lem:cumulative_mono}
    \lean{MPSTensor.cumulativeSpan_mono}
    \leanok
    \uses{def:cumulative_span}
    $T_n(A) \le T_{n+1}(A)$ for all~$n$.
\end{lemma}

\begin{proof}\leanok\uses{def:cumulative_span}
    $T_{n+1}$ includes all words of length $\le n+1$,
    which contains all words of length $\le n$.
\end{proof}

\begin{lemma}[Dimension bound]\label{lem:cumulative_finrank_le}
    \lean{MPSTensor.cumulativeSpan_finrank_le}
    \leanok
    \uses{def:cumulative_span}
    $\dim T_n(A) \le D^2$ for all~$n$.
\end{lemma}

\begin{proof}\leanok
    $T_n(A) \subseteq \MN{D}$ and $\dim \MN{D} = D^2$.
\end{proof}

\begin{theorem}[Cumulative span stabilisation]\label{thm:cumulative_stable}
    \lean{MPSTensor.cumulativeSpan_stable}
    \leanok
    \uses{def:cumulative_span}
    If $T_n(A) = T_{n+1}(A)$, then $T_m(A) = T_n(A)$
    for all $m \ge n$.
\end{theorem}

\begin{proof}
    \leanok
    If $T_n = T_{n+1}$, then $S_{n+1} \subseteq T_n$.
    Left-multiplying any element of $S_{n+1}$ by $A^i$ gives
    an element of $S_{n+2}$, so $S_{n+2} \subseteq T_{n+1} = T_n$.
    By induction, $S_m \subseteq T_n$ for all $m > n$,
    hence $T_m = T_n$.
\end{proof}

\begin{theorem}[Dimension growth]\label{thm:cumulative_strict_mono}
    \lean{MPSTensor.cumulativeSpan_finrank_strict_mono}
    \leanok
    \uses{def:cumulative_span}
    If $T_n(A) \neq T_{n+1}(A)$, then
    $\dim T_{n+1}(A) > \dim T_n(A)$.
\end{theorem}

\begin{proof}
    \leanok
    \uses{lem:cumulative_mono}
    $T_n \le T_{n+1}$ (Lemma~\ref{lem:cumulative_mono})
    and $T_n \neq T_{n+1}$
    imply strict submodule inclusion, hence strict rank increase.
\end{proof}

\begin{lemma}[Word span characterises block injectivity]\label{lem:wordspan_blk_inj}
    \lean{MPSTensor.wordSpan_eq_top_iff_isNBlkInjective}
    \leanok
    \uses{def:word_span, def:l_blk_injective}
    $S_L(A) = \MN{D}$ if and only if $A$ is $L$-block injective.
\end{lemma}

\begin{proof}\leanok\uses{def:word_span, def:l_blk_injective}
    Both conditions assert that the span of all products of
    length~$L$ equals $\MN{D}$.
\end{proof}

\section{Fitting decomposition}

\begin{definition}[Fitting decomposition]\label{def:fitting_decomp}
    \lean{Wielandt.FittingDecomposition}
    \leanok
    A \emph{Fitting decomposition} of a linear endomorphism
    $f : V \to V$ on a finite-dimensional vector space over an
    algebraically closed field consists of:
    \begin{enumerate}
        \item $f$ is nilpotent on the generalized $0$-eigenspace,
        \item $f$ is invertible on each generalized $\mu$-eigenspace
              for $\mu \neq 0$,
        \item the generalized eigenspaces span~$V$,
        \item the generalized eigenspaces are linearly independent.
    \end{enumerate}
\end{definition}

\begin{theorem}[Fitting decomposition exists]\label{thm:fitting_decomp}
    \lean{Wielandt.fittingDecomposition}
    \leanok
    \uses{def:fitting_decomp}
    Every linear endomorphism on a finite-dimensional vector space
    over an algebraically closed field admits a Fitting decomposition.
\end{theorem}

\begin{proof}
    \leanok
    Nilpotency on the zero generalized eigenspace follows from
    the definition of generalized eigenspaces.
    Invertibility on nonzero generalized eigenspaces follows
    because $f - \mu$ is nilpotent there, so $f = \mu(1 - (1-f/\mu))$
    is invertible.
    Spanning and independence are standard results for
    generalized eigenspaces over algebraically closed fields.
    In~\cite{Sanz2010Quantum} and \cite[Lemma~6.3(b)]{Wolf2012Quantum},
    the Jordan normal form is used directly.
    Our proof uses generalized eigenspaces as a substitute
    for the full Jordan normal form.
\end{proof}

\begin{theorem}[Nilpotent power bound]\label{thm:nilpotent_pow_bound}
    \lean{Wielandt.nilpotent_pow_eq_zero_of_finrank}
    \leanok
    If $f$ is nilpotent on a space of dimension~$n$,
    then $f^n = 0$.
\end{theorem}

\begin{proof}\leanok
    A nilpotent endomorphism on an $n$-dimensional space has
    nilpotency index $\le n$.
\end{proof}

\section{Nonzero trace product}

\begin{theorem}[Cumulative span reaches $\MN{D}$ --- explicit bound]\label{thm:cumulative_eq_top_bound}
    \lean{MPSTensor.cumulativeSpan_eq_top_of_isNormal_bound}
    \leanok
    \uses{def:cumulative_span, def:normal,
          thm:cumulative_stable, thm:cumulative_strict_mono,
          lem:cumulative_finrank_le}
    If $A$ is normal, then there exists $n \le D^2$ with
    $T_n(A) = \MN{D}$.
\end{theorem}

\begin{proof}
    \leanok
    By Theorem~\ref{thm:cumulative_strict_mono},
    if $T_n \neq T_{n+1}$ then $\dim T_{n+1} > \dim T_n$.
    Since $\dim T_n \le D^2$ (Lemma~\ref{lem:cumulative_finrank_le}),
    after at most $D^2$ steps either $T_n = \MN{D}$ or
    $T_n$ has stabilised.
    Stabilisation before reaching $\MN{D}$ contradicts normality
    (which requires $S_{L_0} = \MN{D}$ for some $L_0$,
    hence $T_{L_0} = \MN{D}$).
\end{proof}

\begin{theorem}[Cumulative span reaches $\MN{D}$]\label{thm:cumulative_eq_top}
    \lean{MPSTensor.cumulativeSpan_eq_top}
    \leanok
    \uses{thm:cumulative_eq_top_bound}
    If $A$ is normal, then $T_{D^2}(A) = \MN{D}$.
\end{theorem}

\begin{proof}\leanok\uses{thm:cumulative_eq_top_bound}
    Take $n$ from Theorem~\ref{thm:cumulative_eq_top_bound};
    since $n \le D^2$ and $T$ is monotone, $T_{D^2} = \MN{D}$.
\end{proof}

\begin{theorem}[Nonzero trace product]\label{thm:nonzero_trace_word}
    \lean{MPSTensor.exists_nonzero_trace_word}
    \leanok
    \uses{thm:cumulative_eq_top}
    If $A$ is normal, then there exists a word $w$ of length
    $\le D^2$ such that $\tr(A^w) \neq 0$.
    This is \cite[Lemma~1]{Sanz2010Quantum}.
\end{theorem}

\begin{proof}
    \leanok
    \uses{thm:cumulative_eq_top}
    By Theorem~\ref{thm:cumulative_eq_top},
    $\Id \in T_{D^2}(A)$,
    so $\Id$ is a linear combination of word products
    of length $\le D^2$.
    At least one has nonzero trace
    (since $\tr(\Id) = D \neq 0$).
\end{proof}

\section{Eigenvector extraction}

\begin{theorem}[Nonzero trace implies eigenvector]\label{thm:eigenvector_from_trace}
    \lean{exists_eigenvector_of_trace_ne_zero}
    \leanok
    \uses{def:fitting_decomp}
    If $M \in \MN{D}$ has $\tr(M) \neq 0$, then $M$ has a
    nonzero eigenvalue $\mu \neq 0$ and a corresponding eigenvector
    $\varphi \neq 0$ with $M \varphi = \mu \varphi$.
\end{theorem}

\begin{proof}
    \leanok
    \uses{thm:fitting_decomp}
    A matrix with nonzero trace has a nonzero eigenvalue
    (since trace = sum of eigenvalues counted with multiplicity).
    The Fitting decomposition
    (Theorem~\ref{thm:fitting_decomp})
    ensures the corresponding generalized eigenspace is nontrivial,
    yielding an eigenvector.
\end{proof}

\begin{theorem}[Eigenvalue extraction]\label{thm:eigenvalue_extraction}
    \lean{MPSTensor.exists_word_eigenvector}
    \leanok
    \uses{thm:nonzero_trace_word, thm:eigenvector_from_trace}
    If $A$ is normal, then there exist a word $w_0$ with
    $|w_0| \le D^2$,
    a nonzero scalar $\mu \neq 0$, and a nonzero vector
    $\varphi \neq 0$ such that $A^{w_0} \varphi = \mu \varphi$.
\end{theorem}

\begin{proof}
    \leanok
    \uses{thm:nonzero_trace_word, thm:eigenvector_from_trace}
    By Theorem~\ref{thm:nonzero_trace_word}, there is a word $w_0$
    with $\tr(A^{w_0}) \neq 0$ and $|w_0| \le D^2$.
    By Theorem~\ref{thm:eigenvector_from_trace},
    $A^{w_0}$ has a nonzero eigenvector.
\end{proof}

\section{Eigenvector spreading}

\begin{definition}[Cumulative vector span]\label{def:cumulative_vector_span}
    \lean{MPSTensor.cumulativeVectorSpan}
    \leanok
    \uses{def:mps_tensor}
    Given an MPS tensor~$A$ and a vector
    $\varphi \in \C^D$, the \emph{cumulative vector span} is
    \[
        K_n(A, \varphi) = 
        \spn_{\C}\{A^w \varphi : |w| \le n\}
        \subseteq  \C^D.
    \]
    This is \cite[proof of Lemma~2(a)]{Sanz2010Quantum}.
\end{definition}

\begin{theorem}[Vector span stabilisation]\label{thm:vector_span_stable}
    \lean{MPSTensor.cumulativeVectorSpan_stable}
    \leanok
    \uses{def:cumulative_vector_span}
    If $K_n(A, \varphi) = K_{n+1}(A, \varphi)$, then
    $K_m(A, \varphi) = K_n(A, \varphi)$ for all $m \ge n$.
\end{theorem}

\begin{proof}\leanok
    Same argument as Theorem~\ref{thm:cumulative_stable}:
    left-multiplication by $A^i$ cannot escape the stabilised span.
\end{proof}

\begin{lemma}[Vector span from matrix span]\label{lem:vector_span_from_matrix}
    \lean{MPSTensor.cumulativeVectorSpan_eq_top_of_cumulativeSpan_eq_top}
    \leanok
    \uses{def:cumulative_vector_span, def:cumulative_span}
    If $T_N(A) = \MN{D}$ and $\varphi \neq 0$, then
    $K_N(A, \varphi) = \C^D$.
\end{lemma}

\begin{proof}
    \leanok
    Every matrix $M \in \MN{D}$ is in $T_N(A)$, so in particular
    every rank-one matrix sending $\varphi$ to any target vector~$v$
    is in $T_N$.
    Applying such matrices to $\varphi$ recovers all of $\C^D$.
\end{proof}

\begin{theorem}[Eigenvector spreading]\label{thm:eigenvector_spreading}
    \lean{MPSTensor.eigenvector_spreading}
    \leanok
    \uses{def:cumulative_vector_span, def:normal}
    Let $A$ be a normal MPS tensor and let
    $A^{i_0} \varphi = \mu \varphi$ with $\mu \neq 0$
    and $\varphi \neq 0$.
    Then $K_{D-1}(A, \varphi) = \C^D$.
    This is \cite[Lemma~2(a)]{Sanz2010Quantum}.
\end{theorem}

\begin{proof}
    \leanok
    \uses{thm:cumulative_eq_top, lem:vector_span_from_matrix,
          thm:vector_span_stable}
    The vector span $K_n$ is monotone and $\dim K_n \le D$.
    If $K_n = K_{n+1}$ for some $n < D-1$, the span stabilises
    (Theorem~\ref{thm:vector_span_stable}).
    But $T_{D^2}(A) = \MN{D}$ by Theorem~\ref{thm:cumulative_eq_top},
    so Lemma~\ref{lem:vector_span_from_matrix} gives
    $K_{D^2}(A, \varphi) = \C^D$, contradicting early stabilisation.
    Hence $\dim K_n$ grows by $\ge 1$ at each step,
    reaching $D$ by step $D - 1$.
\end{proof}

\section{Assembly}

\begin{definition}[Wielandt analysis]\label{def:wielandt_analysis}
    \lean{MPSTensor.WielandtAnalysis}
    \leanok
    \uses{def:mps_tensor, def:normal}
    A \emph{Wielandt analysis} for a normal MPS tensor~$A$
    packages:
    a word $w_0$ of length $\le D^2$,
    a nonzero eigenvalue $\mu$,
    and a nonzero eigenvector $\varphi$
    of the matrix $A^{w_0}$, such that
    $A^{w_0} \varphi = \mu \varphi$.
\end{definition}

\begin{theorem}[Wielandt analysis exists]\label{thm:wielandt_analysis}
    \lean{MPSTensor.wielandtAnalysis}
    \leanok
    \uses{def:wielandt_analysis, thm:eigenvalue_extraction}
    Every normal MPS tensor admits a Wielandt analysis.
\end{theorem}

\begin{proof}
    \leanok
    \uses{thm:eigenvalue_extraction}
    Immediate from Theorem~\ref{thm:eigenvalue_extraction}.
\end{proof}

\begin{theorem}[Wielandt chain]\label{thm:wielandt_chain}
    \lean{MPSTensor.wielandt_chain}
    \leanok
    \uses{thm:wielandt_analysis, thm:eigenvector_spreading}
    Every normal MPS tensor admits a Wielandt analysis
    (Theorem~\ref{thm:wielandt_analysis}), and the eigenvector
    from that analysis spreads to all of $\C^D$ in $D-1$ steps
    (Theorem~\ref{thm:eigenvector_spreading}).
\end{theorem}

\begin{proof}
    \leanok
    \uses{thm:wielandt_analysis, thm:eigenvector_spreading}
    Combine Theorems~\ref{thm:wielandt_analysis}
    and~\ref{thm:eigenvector_spreading}.
\end{proof}

\begin{theorem}[Wielandt bound]\label{thm:wielandt_bound}
    \lean{MPSTensor.isNormal_implies_cumulativeSpan}
    \leanok
    \uses{def:cumulative_span, def:normal}
    If $A$ is a normal MPS tensor with bond dimension~$D$,
    then $T_{D^2}(A) = \MN{D}$.
    In particular, the cumulative span stabilises at index
    at most $D^2$.
    This is \cite[Theorem~1]{Sanz2010Quantum}
    (which gives the sharper bound $D^2 - d + 1$ in certain
    cases; here we prove the $D^2$ bound).
\end{theorem}

\begin{proof}
    \leanok
    \uses{thm:cumulative_eq_top}
    Immediate from Theorem~\ref{thm:cumulative_eq_top}.
\end{proof}

\section{Lemma~2(b): fixed-length matrix spanning}\label{sec:lemma2b}

The preceding sections establish Lemma~2(a) of~\cite{Sanz2010Quantum}:
eigenvector spreading gives a cumulative vector span $K_{D-1}(A,\varphi) = \C^D$.
The next step, Lemma~2(b), strengthens this to \emph{fixed-length}
matrix spanning:
there exists an $n_0$ such that $S_{n_0}(A) = \MN{D}$
(i.e.\ word products of \emph{exactly} one length span all matrices).
The following results cover the
algebraic assembly part of this argument.

\begin{definition}[Fixed-length vector span]\label{def:vector_spread_span}
    \lean{MPSTensor.vectorSpreadSpan}
    \leanok
    \uses{def:mps_tensor, def:eval_word}
    The \emph{fixed-length vector span} at length~$n$ is
    \[
        H_n(A, \varphi) = \spn_{\C}\{A^w \varphi : |w| = n\}
        \subseteq \C^D.
    \]
    This is $S_n(A)\ket{\varphi}$ in the notation
    of~\cite{Sanz2010Quantum}.
\end{definition}

\begin{lemma}[Word span products]\label{lem:word_span_mul}
    \lean{MPSTensor.wordSpan_mul_le}
    \leanok
    \uses{def:word_span}
    $S_m(A) \cdot S_n(A) \subseteq S_{m+n}(A)$:
    the product of word spans at lengths $m$ and $n$
    is contained in the word span at length $m + n$.
\end{lemma}

\begin{proof}\leanok
    Every product $A^{w_1} A^{w_2}$ with $|w_1| = m$
    and $|w_2| = n$ equals $A^{w_1 w_2}$ with $|w_1 w_2| = m + n$.
\end{proof}

\begin{lemma}[Eigenvector padding]\label{lem:eigenvector_padding}
    \lean{MPSTensor.cumulativeVectorSpan_eq_vectorSpreadSpan_of_eigenvector}
    \leanok
    \uses{def:cumulative_vector_span, def:vector_spread_span}
    Suppose $A^{i_0} \varphi = \mu \varphi$ with $\mu \neq 0$.
    Then for all~$n$,
    \[
        K_n(A, \varphi) = H_n(A, \varphi).
    \]
    In particular, if $K_n(A,\varphi) = \C^D$ then
    already $H_n(A,\varphi) = \C^D$.
\end{lemma}

\begin{proof}
    \leanok
    Any word $w$ with $|w| \le n$ can be padded to exact
    length~$n$ by appending $n - |w|$ copies of the index~$i_0$.
    The eigenvector relation gives
    $A^{w \cdot i_0^k} \varphi = \mu^k \cdot A^w \varphi$,
    so the padded vector is a nonzero scalar multiple
    of the original.
    Hence every generator of $K_n$ lies in $H_n$.
    The reverse inclusion $H_n \le K_n$ is immediate.
\end{proof}

\begin{theorem}[Vector spanning to matrix spanning (assembly step)]%
    \label{thm:lemma2b_assembly}
    \lean{MPSTensor.wordSpan_eq_top_of_vectorSpreadSpan_eq_top_of_rankOneBasis}
    \leanok
    \uses{def:word_span, def:vector_spread_span, lem:word_span_mul}
    Assume:
    \begin{enumerate}
        \item $H_n(A, \varphi) = \C^D$
              (length-$n$ word products applied to $\varphi$
              span all of~$\C^D$), and
        \item for each standard basis vector $e_j$,
              the rank-one operator
              $\ketbra{\varphi}{e_j}$
              lies in $S_m(A)$.
    \end{enumerate}
    Then $S_{n+m}(A) = \MN{D}$.
\end{theorem}

\begin{proof}
    \leanok
    \uses{lem:word_span_mul}
    Fix a matrix unit $E_{ij}$.
    By hypothesis~(1), there exists $M_i \in S_n(A)$
    with $M_i \varphi = e_i$.
    By hypothesis~(2), $\ketbra{\varphi}{e_j} \in S_m(A)$.
    Their product satisfies
    \[
        M_i \cdot \ketbra{\varphi}{e_j}
        = \ketbra{M_i \varphi}{e_j}
        = \ketbra{e_i}{e_j}
        = E_{ij},
    \]
    and $E_{ij} \in S_{n+m}(A)$ by
    Lemma~\ref{lem:word_span_mul}.
    Since the matrix units span $\MN{D}$,
    $S_{n+m}(A) = \MN{D}$.
\end{proof}

\begin{lemma}[Blocking transfer]\label{lem:blocking_transfer}
    \lean{MPSTensor.wordSpan_eq_top_of_blockTensor_wordSpan_eq_top}
    \leanok
    \uses{def:word_span, def:blocked_tensor}
    If the blocked tensor $A^{[L]}$ satisfies
    $S_n(A^{[L]}) = M_{D}(\C)$,
    then $S_{nL}(A) = M_{D}(\C)$.
\end{lemma}

\begin{proof}
    \leanok
    Each length-$n$ word in the blocked alphabet
    $\{0,\ldots,d^L-1\}$ decodes to a length-$nL$ word
    in the original alphabet $\{0,\ldots,d-1\}$,
    so $S_n(A^{[L]}) \subseteq S_{nL}(A)$.
\end{proof}

\begin{remark}[Status of Lemma~2(b)]\label{rem:lemma2b_status}
    Theorem~\ref{thm:lemma2b_assembly} establishes the
    \emph{assembly} part of Lemma~2(b)
    of~\cite{Sanz2010Quantum}:
    rank-one operators plus fixed-length vector spanning
    imply fixed-length matrix spanning.
    The remaining \emph{rank-one construction step}
    is completed in \S\ref{sec:rectangular_span}:
    Theorem~\ref{thm:rankone_extraction_blockTensor}
    produces the required rank-one element via the
    Fitting decomposition, and
    Theorem~\ref{thm:wielandt_lemma2b} gives the fully
    unconditional Lemma~2(b).
\end{remark}

\section{Blocked tensor and rectangular span}\label{sec:rectangular_span}

This section extends the Lemma~2(b) infrastructure to the
\emph{blocked tensor}~$A^{[L]}$ and develops the ``rectangular span''
used to produce rank-one elements.

\begin{theorem}[Blocking preserves normality]\label{thm:isNormal_blockTensor}
    \lean{MPSTensor.isNormal_blockTensor}
    \leanok
    \uses{def:normal, def:blocked_tensor, def:word_span}
    If $A$ is normal and $L > 0$, then the blocked tensor
    $A^{[L]}$ is also normal.
\end{theorem}

\begin{proof}\leanok
    \uses{lem:word_span_le_blockTensor}
    If $S_N(A) = \MN{D}$ then
    $S_{NL}(A) = \MN{D}$
    (padding by repeated single-letter words).
    By Lemma~\ref{lem:word_span_le_blockTensor},
    $S_{NL}(A) \subseteq S_N(A^{[L]})$,
    so $S_N(A^{[L]}) = \MN{D}$.
\end{proof}

\begin{lemma}[Chunking: word span embeds into blocked word span]%
    \label{lem:word_span_le_blockTensor}
    \lean{MPSTensor.wordSpan_le_wordSpan_blockTensor}
    \leanok
    \uses{def:word_span, def:blocked_tensor}
    $S_{nL}(A) \subseteq S_n(A^{[L]})$ for all $n, L$.
\end{lemma}

\begin{proof}\leanok
    \uses{lem:word_span_mul}
    Every length-$nL$ word in the original alphabet
    can be split into $n$ consecutive chunks of length~$L$,
    each of which is a single blocked Kraus operator.
    By induction on~$n$, using the word span product lemma
    (Lemma~\ref{lem:word_span_mul}).
\end{proof}

\begin{theorem}[Word eigenvector gives blocked eigenvector]%
    \label{thm:blockTensor_single_eigenvector}
    \lean{MPSTensor.blockTensor_single_eigenvector}
    \leanok
    \uses{def:blocked_tensor}
    If $A^{w_0} \varphi = \mu \varphi$ for a word $w_0$
    of length~$L$, then
    $(A^{[L]})^{j_0} \varphi = \mu \varphi$,
    where $j_0$ is the encoded blocked index corresponding
    to $w_0$.
\end{theorem}

\begin{proof}\leanok
    By definition, $(A^{[L]})^{j_0} = A^{w_0}$.
\end{proof}

\begin{theorem}[Wielandt blocked assembly]\label{thm:wielandt_blocked_assembly}
    \lean{MPSTensor.wielandt_blocked_assembly}
    \leanok
    \uses{def:word_span, def:normal, def:blocked_tensor,
          thm:isNormal_blockTensor, thm:blockTensor_single_eigenvector,
          thm:eigenvector_spreading, lem:blocking_transfer}
    Suppose $A$ is normal and $L > 0$, and we have:
    \begin{enumerate}
        \item a word $w_0$ of length $L$ with eigenvector
              $A^{w_0} \varphi = \mu \varphi$
              ($\mu \neq 0$, $\varphi \neq 0$),
        \item a word $\tau_0$ of length $L$ with transpose
              eigenvector
              $(A^{\tau_0})^T \psi = \nu \psi$
              ($\nu \neq 0$, $\psi \neq 0$),
        \item a rank-one element
              $\varphi \psi^T \in S_m(A^{[L]})$
              for some $m$.
    \end{enumerate}
    Then $S_{((D-1) + m + (D-1)) \cdot L}(A) = \MN{D}$.
\end{theorem}

\begin{proof}
    \leanok
    \uses{thm:isNormal_blockTensor, thm:blockTensor_single_eigenvector,
          thm:eigenvector_spreading, lem:blocking_transfer,
          thm:lemma2b_assembly}
    Apply the conditional Lemma~2(b) assembly
    (Theorem~\ref{thm:lemma2b_assembly}) to the
    blocked tensor $A^{[L]}$,
    which is normal by Theorem~\ref{thm:isNormal_blockTensor}.
    The eigenvector and transpose eigenvector are transferred
    to single-index eigenvectors of $A^{[L]}$ by
    Theorem~\ref{thm:blockTensor_single_eigenvector}.
    Transfer back to the original tensor via
    Lemma~\ref{lem:blocking_transfer}.
\end{proof}

\begin{theorem}[Rank-one extraction for blocked tensor]%
    \label{thm:rankone_extraction_blockTensor}
    \lean{MPSTensor.exists_rankOne_mem_wordSpan_blockTensor}
    \leanok
    \uses{def:word_span, def:normal, def:blocked_tensor,
          thm:isNormal_blockTensor, thm:nonzero_trace_word,
          thm:eigenvector_from_trace}
    Suppose $A$ is normal and $N_0 \ge 1$ with
    $S_{N_0}(A) = \MN{D}$.
    Then there exist words $\sigma_0, \tau_0$ of length~$N_0$,
    nonzero vectors $\varphi, \psi$, nonzero scalars $\mu, \nu$,
    and $m \in \N$ such that:
    \begin{enumerate}
        \item $A^{\sigma_0} \varphi = \mu \varphi$,
        \item $(A^{\tau_0})^T \psi = \nu \psi$,
        \item $\varphi \psi^T \in S_m(A^{[N_0]})$.
    \end{enumerate}
\end{theorem}

\begin{proof}
    \leanok
    \uses{thm:isNormal_blockTensor, thm:nonzero_trace_word,
          thm:eigenvector_from_trace}
    Set $B = A^{[N_0]}$.
    Since $S_{N_0}(A) = \MN{D}$,
    a nonzero-trace word exists
    (Theorem~\ref{thm:nonzero_trace_word}),
    yielding a column eigenvector $\varphi$
    via Theorem~\ref{thm:eigenvector_from_trace}.
    A row eigenvector $\psi$ is obtained analogously
    from the transpose tensor.
    Set $P = (B^{j_0})^D$ and $Q = (B^{k_0})^D$,
    where $j_0, k_0$ encode $\sigma_0, \tau_0$
    in the blocked alphabet.
    The Fitting decomposition ensures
    $\varphi \in \operatorname{range}(P)$
    and $\psi \in \operatorname{range}(Q^T)$,
    so $\varphi \psi^T \in
    \operatorname{range}(\mathrm{mulLeft}(P)
    \circ \mathrm{mulRight}(Q))
    \subseteq S_{2D}(B)$.
\end{proof}

\begin{theorem}[Wielandt Lemma~2(b)]\label{thm:wielandt_lemma2b}
    \lean{MPSTensor.wielandt_lemma2b}
    \leanok
    \uses{def:word_span, def:normal,
          thm:rankone_extraction_blockTensor,
          thm:wielandt_blocked_assembly}
    If $A$ is a normal MPS tensor with bond dimension~$D$,
    then there exists $N$ such that $S_N(A) = \MN{D}$.
    This is \cite[Lemma~2(b)]{Sanz2010Quantum}.
\end{theorem}

\begin{proof}
    \leanok
    \uses{thm:rankone_extraction_blockTensor,
          thm:wielandt_blocked_assembly}
    By normality, some $N_0$ has $S_{N_0}(A) = \MN{D}$.
    If $N_0 = 0$ we are done.
    Otherwise, Theorem~\ref{thm:rankone_extraction_blockTensor}
    produces eigenvectors and a rank-one element
    $\varphi \psi^T \in S_m(A^{[N_0]})$.
    Theorem~\ref{thm:wielandt_blocked_assembly} then gives
    $S_{((D-1) + m + (D-1)) \cdot N_0}(A) = \MN{D}$.
\end{proof}

\begin{remark}[Gap closed for Lemma~2(b) via blocking]%
    \label{rem:blocked_lemma2b_gap}
    The rank-one extraction gap noted in earlier versions of this
    blueprint has been closed:
    Theorem~\ref{thm:rankone_extraction_blockTensor} produces the
    required rank-one element, and
    Theorem~\ref{thm:wielandt_lemma2b} gives the fully unconditional
    Lemma~2(b).
    See Theorem~\ref{thm:rankone_extraction_blockTensor} and
    Theorem~\ref{thm:wielandt_lemma2b} above.
\end{remark}

\section{Primitivity and normality}

\begin{theorem}[Primitive MPS tensors have nonzero word products]\label{thm:prim_nonzero_word}
    \lean{MPSTensor.exists_nonzero_evalWord_of_isPrimitive}
    \leanok
    \uses{def:eval_word}
    If $A$ is a primitive MPS tensor in DS gauge,
    then for every $n \ge 0$ there exists a word $w$ of length
    \emph{exactly}~$n$ with $A^w \neq 0$.
\end{theorem}

\begin{proof}\leanok
    The PSD fixed point $\rho$ of $\E_A$ satisfies
    $\E_A^n(\rho) = \rho \neq 0$.
    Since $\E_A^n(\rho) = \sum_{|w|=n} A^w \rho (A^w)^\dagger$,
    at least one summand is nonzero, hence some $A^w \neq 0$.
\end{proof}

\begin{theorem}[Iterated transfer does not vanish]\label{thm:transfer_pow_ne_zero}
    \lean{MPSTensor.transferMap_pow_ne_zero_of_isPrimitive}
    \leanok
    \uses{def:transfer_map, thm:prim_nonzero_word}
    If $A$ is a primitive MPS tensor (with DS gauge), then
    $\E_A^n \neq 0$ for all $n \ge 0$.
\end{theorem}

\begin{proof}\leanok\uses{thm:prim_nonzero_word}
    By Theorem~\ref{thm:prim_nonzero_word},
    there exists a word~$w$ of length~$n$ with $A^w \neq 0$,
    so $\E_A^n(\rho)$ has a nonzero summand.
\end{proof}

\subsection{Primitivity-to-normality bridge}\label{subsec:prim_to_normal}

The spectral-gap notion of primitivity
(unique PSD fixed point $\rho$ with $\rho_{\mathrm{spec}}(\E_A - P_\rho) < 1$)
is connected to normality through the following results.

\begin{theorem}[Fixed-point uniqueness from spectral gap]%
    \label{thm:fixedpoint_unique_spectral}
    \lean{MPSTensor.IsPrimitiveMPS.fixedPoint_unique}
    \leanok
    If $A$ is a primitive MPS tensor
    with PSD fixed point $\rho$,
    then every fixed point $\sigma$ of $\E_A$ is proportional
    to $\rho$: $\sigma = \frac{\tr(\sigma)}{\tr(\rho)} \rho$.
\end{theorem}

\begin{proof}\leanok
    Set $\sigma' = \sigma - \frac{\tr(\sigma)}{\tr(\rho)} \rho$.
    Then $\tr(\sigma') = 0$ and $(E - P_\rho)(\sigma') = \sigma'$.
    If $\sigma' \neq 0$, then $1$ is an eigenvalue of $E - P_\rho$,
    contradicting the spectral gap
    $\rho_{\mathrm{spec}}(E - P_\rho) < 1$.
\end{proof}

\begin{theorem}[Complement powers tend to zero]%
    \label{thm:complement_pow_zero}
    \lean{MPSTensor.IsPrimitiveMPS.complement_pow_tendsto_zero}
    \leanok
    If $A$ is a primitive MPS tensor with PSD fixed point $\rho$,
    then $(\E_A - P_\rho)^n \to 0$ in operator norm.
\end{theorem}

\begin{proof}\leanok
    The spectral radius of $\E_A - P_\rho$ is strictly less
    than $1$ by the spectral-gap hypothesis.
    Apply the general result that powers of an operator
    with spectral radius $< 1$ tend to zero.
\end{proof}

\begin{remark}[Primitivity: spectral gap vs.\ positive definiteness]%
    \label{rem:prim_definition_mismatch}
    The definition of primitivity used in \S\ref{subsec:prim_to_normal}
    encodes only the spectral-gap condition: there exists a nonzero PSD
    fixed point $\rho$ and $\rho_{\mathrm{spec}}(\E_A - P_\rho) < 1$.
    The notion of ``primitive'' in
    \cite[Proposition~3]{Sanz2010Quantum} additionally requires
    $\rho$ to be positive definite (full rank).
    Without positive definiteness, the spectral-gap condition alone
    does \emph{not} imply irreducibility or normality.
    A counterexample is $D = 2$, $\rho = \diag(1,0)$.

    Both gaps identified in earlier versions of this remark
    are now closed:
    \begin{enumerate}
        \item Primitive $+$ $\rho$ positive definite
              $\Rightarrow$ irreducible tensor: 
              Theorem~\ref{thm:primitive_implies_irreducible}
              (\S\ref{subsec:prim_implies_irred}).
        \item Cumulative spanning $+$ aperiodicity
              $\Rightarrow$ fixed-length spanning (normality):
              Theorem~\ref{thm:cumspan_to_normal}
              (\S\ref{sec:cumspan_to_wordspan}).
    \end{enumerate}
    Together with the Burnside bridge
    (Chapter~\ref{ch:canonical}, \S\ref{subsec:burnside_bridge}),
    the full chain is:
    \[
        \text{primitive} + \rho\text{ PosDef}
        \;\Rightarrow\;
        \text{irreducible}
        \;\Rightarrow\;
        \operatorname{algSpan}(A) = \MN{D}
        \;\Rightarrow\;
        \text{normal}.
    \]
    The last step uses the aperiodicity condition
    $\Id \in S_1(A)$ (Remark~\ref{rem:aperiodicity}).
\end{remark}

\subsection{Primitive implies irreducible}\label{subsec:prim_implies_irred}

\begin{theorem}[Primitive with PosDef fixed point implies irreducible tensor]%
    \label{thm:primitive_implies_irreducible}
    \lean{MPSTensor.isIrreducibleTensor_of_isPrimitiveMPS_of_posDef}
    \leanok
    \uses{def:irreducible_tensor, def:has_inv_proj,
          thm:complement_pow_zero}
    If $A$ is a primitive MPS tensor with PSD fixed point $\rho$
    and $\rho$ is positive definite, then $A$ is an irreducible tensor.
    This is part of \cite[Proposition~3]{Sanz2010Quantum}
    (see also \cite[\S2.3]{Cirac2016Density}).
\end{theorem}

\begin{proof}
    \leanok
    \uses{thm:complement_pow_zero}
    By contrapositive.
    Assume $A$ has an invariant projection~$P$
    ($P \neq 0$, $P \neq \Id$, $(\Id - P) A^i P = 0$ for all~$i$).
    Set $\sigma = P \rho P$.
    \begin{enumerate}
        \item \emph{Invariance of $\sigma$ under $\E_A^n$:}
              the projection relation
              $(\Id - P) A^i P = 0$ implies
              $(\Id - P) \cdot \E_A^n(\sigma) = 0$ for all~$n$
              (by induction on~$n$).
        \item \emph{Convergence:}
              By Theorem~\ref{thm:complement_pow_zero},
              $\E_A^n(\sigma) \to \frac{\tr(\sigma)}{\tr(\rho)} \rho$.
              Hence
              $(\Id - P) \cdot \frac{\tr(\sigma)}{\tr(\rho)} \rho = 0$.
        \item \emph{Nonzero scalar:}
              Since $P \neq 0$ and $\rho$ is positive definite,
              $\tr(P \rho P) > 0$, so
              $\frac{\tr(\sigma)}{\tr(\rho)} \neq 0$.
              Therefore $(\Id - P) \rho = 0$.
        \item \emph{Contradiction:}
              Since $\rho$ is positive definite, it is a unit
              in $\MN{D}$.
              From $(\Id - P) \rho = 0$ we get $\Id - P = 0$,
              i.e.\ $P = \Id$, contradicting $P \neq \Id$.
    \end{enumerate}
\end{proof}

\section{Cumulative span to word span}\label{sec:cumspan_to_wordspan}

Cumulative spanning ($T_N(A) = \MN{D}$, the union of word spans
of all lengths~$\le N$) is weaker than normality
($S_L(A) = \MN{D}$ for a single length~$L$).
The gap is closed by an aperiodicity hypothesis.

\begin{remark}[Aperiodicity condition]\label{rem:aperiodicity}
    We say $A$ satisfies the \emph{aperiodicity condition}
    if $\Id \in S_1(A)$, i.e.\ the identity matrix lies in the
    span of the Kraus operators $\{A^i\}$.
    For a doubly-stochastic tensor, $\sum_i A^i (A^i)^\dagger = \Id$
    and $\sum_i (A^i)^\dagger A^i = \Id$, which gives
    $\Id \in S_1(A)$.
    More generally, blocking to a primitive transfer map
    (Theorem~\ref{thm:blocking_primitivity}) yields a tensor
    with $\Id \in S_1(A^{[p]})$.
\end{remark}

\begin{remark}[Counterexample: cumulative spanning does not imply normality]%
    \label{rem:cumspan_counterexample}
    In $\MN{2}$, let $A^0 = E_{12}$ and $A^1 = E_{21}$
    (the off-diagonal matrix units).
    Then $\operatorname{algSpan}(A) = \MN{2}$ and $T_2(A) = \MN{2}$,
    but the word spans alternate:
    $S_n(A)$ contains only off-diagonal matrices for odd~$n$
    and only diagonal matrices for even~$n$.
    Hence $S_n(A) \neq \MN{2}$ for every~$n$,
    and $A$ is not normal.
    The aperiodicity condition $\Id \in S_1(A)$ fails:
    $S_1(A) = \spn_{\C}\{E_{12}, E_{21}\}$.
\end{remark}

\begin{theorem}[Word span monotonicity from aperiodicity]%
    \label{thm:wordspan_mono_aperiodic}
    \lean{MPSTensor.wordSpan_mono_succ_of_one_mem_wordSpan_one}
    \leanok
    \uses{def:word_span, lem:word_span_mul}
    If $\Id \in S_1(A)$, then
    $S_n(A) \le S_{n+1}(A)$ for all~$n$.
\end{theorem}

\begin{proof}\leanok\uses{lem:word_span_mul}
    For any $M \in S_n(A)$,
    $M = M \cdot \Id \in S_n(A) \cdot S_1(A) \subseteq S_{n+1}(A)$
    by Lemma~\ref{lem:word_span_mul}.
\end{proof}

\begin{theorem}[Cumulative span equals word span under aperiodicity]%
    \label{thm:cumspan_eq_wordspan}
    \lean{MPSTensor.cumulativeSpan_eq_wordSpan_of_one_mem_wordSpan_one}
    \leanok
    \uses{def:cumulative_span, def:word_span,
          thm:wordspan_mono_aperiodic}
    If $\Id \in S_1(A)$, then
    $T_n(A) = S_n(A)$ for all~$n$.
\end{theorem}

\begin{proof}\leanok\uses{thm:wordspan_mono_aperiodic}
    By Theorem~\ref{thm:wordspan_mono_aperiodic},
    the word spans are monotone: $S_0 \le S_1 \le \cdots$.
    The supremum $T_n = S_0 + \cdots + S_n$ therefore equals the
    largest term~$S_n$.
\end{proof}

\begin{theorem}[Cumulative spanning plus aperiodicity implies normality]%
    \label{thm:cumspan_to_normal}
    \lean{MPSTensor.isNormal_of_cumulativeSpan_eq_top_of_aperiodic}
    \leanok
    \uses{def:normal, def:cumulative_span,
          thm:cumspan_eq_wordspan}
    If $T_N(A) = \MN{D}$ and $\Id \in S_1(A)$, then $A$ is normal.
\end{theorem}

\begin{proof}\leanok\uses{thm:cumspan_eq_wordspan}
    By Theorem~\ref{thm:cumspan_eq_wordspan},
    $S_N(A) = T_N(A) = \MN{D}$.
\end{proof}

\begin{theorem}[Algebra span plus aperiodicity implies normality]%
    \label{thm:algspan_to_normal}
    \lean{MPSTensor.isNormal_of_algSpan_eq_top_of_aperiodic}
    \leanok
    \uses{def:normal, def:alg_span,
          thm:cumspan_to_normal}
    If $\operatorname{algSpan}(A) = \MN{D}$ and $\Id \in S_1(A)$,
    then $A$ is normal.
\end{theorem}

\begin{proof}\leanok\uses{thm:cumspan_to_normal}
    Since $\operatorname{algSpan}(A) = \MN{D}$,
    there exists $N$ with $T_N(A) = \MN{D}$
    (by the ascending chain condition on finite-dimensional subspaces).
    Apply Theorem~\ref{thm:cumspan_to_normal}.
\end{proof}

\begin{theorem}[Irreducible tensor plus aperiodicity implies normality]%
    \label{thm:irred_tensor_aperiodic_normal}
    \lean{MPSTensor.isNormal_of_isIrreducibleTensor_of_aperiodic}
    \leanok
    \uses{def:normal, def:irreducible_tensor,
          thm:algspan_to_normal, thm:burnside_matrix,
          thm:irreducible_tensor_to_action}
    If $A$ is an irreducible tensor, $D > 0$,
    and $\Id \in S_1(A)$, then $A$ is normal.
    This assembles the full chain:
    \[
        \text{irreducible tensor}
        \;\xrightarrow{\ref{thm:irreducible_tensor_to_action}}\;
        \text{irreducible action}
        \;\xrightarrow{\text{Burnside}}\;
        \operatorname{algSpan} = \MN{D}
        \;\xrightarrow{\ref{thm:algspan_to_normal}}\;
        \text{normal}.
    \]
    This corresponds to \cite[Theorem~1 / Proposition~3]{Sanz2010Quantum}.
\end{theorem}

\begin{proof}\leanok
    \uses{thm:algspan_to_normal, thm:burnside_matrix,
          thm:irreducible_tensor_to_action}
    By Theorem~\ref{thm:irreducible_tensor_to_action},
    $A$ acts irreducibly on $\C^D$.
    By Theorem~\ref{thm:burnside_matrix},
    $\operatorname{algSpan}(A) = \MN{D}$.
    By Theorem~\ref{thm:algspan_to_normal} with the
    aperiodicity hypothesis, $A$ is normal.
\end{proof}
